% Paper 2, Section 2 — Related work.
% Literature analysis lifted from talk/imam_ssm_memo_v3.md sections 1-3 (commit fb63e23); the
% SPikE-SSM and LMU papers were read in full, CIM-SSM is paywalled and its numbers are NOT quoted.
% No experimental figure in this section is re-derived; the two of ours that appear
% (spikeout state rate 1.0000; the char-LM tradeoff) are copied from a134f05 / bb79f7f.

\section{Related work}
\label{sec:related}

\subsection{State-space models}

Structured SSMs~\citep{s4,s4d} replaced the nonlinear recurrence of an RNN with a linear one
whose state transition is diagonal (after reparameterization), which is what makes them both
trainable at length and, for our purposes, physically realizable: a diagonal linear recurrence is
a bank of independent leaky integrators. We use an S4D-real--style diagonal decay as the shared
backbone of all four variants, which is the simplest form that keeps the comparison about the
datapath rather than about the recurrence's expressivity.

The design predates the current line. The Legendre Memory Unit~\citep{lmu} derives an optimal
linear delay system from a Pad\'e approximant and projects it onto a Legendre basis, arriving at
a linear SSM independently of the HiPPO/S4 lineage. It matters here for a reason beyond priority:
the LMU was implemented through the Neural Engineering Framework on spiking neuron populations
and \emph{deployed on both Loihi and Braindrop} --- a digital spiking chip and an analog
mixed-signal one --- reaching $97.15\%$ on psMNIST against an LSTM's $89.86\%$. Running an SSM on
neuromorphic hardware, including in a non-spiking analog form, is therefore an existence proof
rather than a proposal. The open question, and ours, is what each route costs.

\subsection{The three datapath routes}

\paragraph{(a) Digital SSMs on digital fabric.} The reference point: run the exact linear
recurrence in a digital datapath, optionally quantized. Our \texttt{digital} variant is this
route at full precision, and it is the baseline every margin in this paper is measured against.

\paragraph{(b) Continuous state, spiking output.} The current ``spiking SSM'' line keeps the
linear state at full precision and applies a spiking nonlinearity only to what leaves the
recurrence. SPikE-SSM~\citep{spikessm} is the clearest instance: a refractory-LIF nonlinearity on
an S4D backbone, with a continuous state $h_t$. Its reported sparsity figures (around $8\%$ on
LRA, $24.5\%$ on WikiText) are therefore \emph{output} sparsity, and its cost is a recall gap
(WikiText perplexity $33.2$ against S4's $21.0$). This distinction is not a technicality; it is
the reason a firing-floor argument about state-carrying spikes does not apply to this route at
all. Our \texttt{spikeout} variant reproduces the design point, and reproduces the signature: its
measured recurrent-state activity is exactly $1.0000$ in every cell we ran, on both tasks and at
every memory load --- all of its sparsity is in the output, none of it in the state.

\paragraph{(c) Native analog state.} Here the device physics \emph{is} the recurrence. CIM-SSM
\citep{cimssm} implements a non-spiking continuous diagonal SSM in a memristor crossbar, using
the device's native short-term-memory relaxation as the state decay --- the exponential
relaxation of the device is $e^{At}$, so there is no discretization mismatch to manage, and
$Ax + Bu$ is an in-memory matrix--vector multiply. Related analog compute-in-memory work
approaches the same substrate from adjacent angles: QS4D~\citep{qs4d} trains diagonal S4D with
quantization-aware training for deployment on tantalum-oxide memristive crossbars and reports
that the training itself confers robustness to analog noise. IMSSA~\citep{imssa} is the same
group's earlier hardware demonstration: recurrent S4D kernels with $A$, $B$ and $C$ programmed
into a single $64\times64$ memristive crossbar and executed step-by-step rather than as a
convolution, with weights quantized to ternary; on a two-class spoken-digit task the deployed
model reaches $81.7\%$ against a $95.1\%$ software reference, a gap the authors attribute to
stuck devices --- an on-silicon instance of analog imperfection charging quality. HPD~\citep{hpd} addresses the
robustness side directly, though in simulation rather than on silicon: a layer-wise
vulnerability analysis of Mamba under weight perturbations --- the noise enters the weights,
not the state --- locates the damage in the final block's output projection, and the remedy is
itself a hybrid datapath split: the projection's SVD factor $U\Sigma$ stays on the analog
array while $V^{\top}$ is offloaded to digital for exact correction, an independent instance
of spending exactness on the side of the datapath the loss depends on. LMU-on-Braindrop is the
earlier precedent. We treat this route in simulation, with the analog storage element's
imperfections modelled explicitly (rails, additive state noise, finite state precision) and
communication by send-on-delta graded events; \S\ref{sec:protocol} states exactly what is and is
not modelled. We do not quote CIM-SSM's reported numbers, which are behind a paywall we could
not clear.

\subsection{What is missing, and what this paper adds}

Two gaps motivate the study.

First, \textbf{the routes are never compared under matched capacity}. Each of (a), (b), (c) is
developed in its own papers against its own baseline, typically a full-precision model of
different size trained under a different schedule, so their efficiency claims are
route-internal and not commensurable. Our four variants share every tensor and every
hyperparameter; the only difference between them is where a nonlinearity sits.

Second, \textbf{energy claims are made at whatever operating point the method happens to reach}.
A variant that emits on $27\%$ of steps and one that emits on $65\%$ are not comparable on quality
without controlling the rate. We therefore sweep the analog event threshold until the analog
variant emits at each spiking comparator's own measured rate, and compare there
(\S\ref{sec:charlm-matched}); on the recall task the three non-digital variants land within
$0.41$--$0.50$ emission natively, so the comparison is already matched.

A third issue is one we created for ourselves and report because it changes published-looking
numbers: \textbf{the comparator's regularization}. The lossy analog state acts as a stochastic
regularizer at small scale, so an analog variant compared against an untuned digital baseline is
credited for a benefit that has nothing to do with its datapath. Giving the baseline the same
training-time state noise removes it, and turns an apparent quality win into an energy-for-quality
tradeoff (\S\ref{sec:charlm-control}). We have not seen this control run in the spiking- or
analog-SSM literature, and the bias it introduces runs in the same direction everywhere: in favour
of the neuromorphic variant.

\subsection{Relation to paper 1}

Paper~1~\citep{paper1} established, on a controlled architecture-fixed protocol, that a recurrent
network carrying its state in spikes cannot sparsify below roughly $50\%$ activity on a
character-LM task, while feed-forward perception sparsifies to $5\%$ and a spiking Transformer to
$2\%$ --- a firing ceiling specific to \emph{recurrence}. It also derived the bound
$\rho \geq H_b^{-1}(\cdot)$ as a candidate explanation. This paper takes up the question that
result raises --- if spikes are a bad way to carry recurrent state, what datapath \emph{should}
carry it? --- and, in \S\ref{sec:bound}, reports that the bound itself does not survive a direct
empirical test at this scale. The empirical ceiling from paper~1 is a measurement and stands; the
bound was a prediction about \emph{why}, and it is the prediction that fails.
